\paragraph{The Tacit Communication Game.} The TCG was developed in Nijmegen to study communication when no shared code exists \citep{deruiter2010}. Because words are ruled out, Senders must design signals for their recipient and Receivers must recognise the intention behind a movement; communicative intent is carried by the parts of a movement that do nothing for the Sender's own goal. Pairs communicate remarkably well, especially with feedback \citep{deruiter2010}. Senders adapt to whom they believe the partner is \citep{newmannorlund2009}, and planning and recognising communicative actions engage a common region (right posterior superior temporal sulcus) whose activity tracks message ambiguity rather than movement complexity \citep{noordzij2009}.

\paragraph{Conceptual alignment.} Stolk and colleagues used the TCG to show that solving novel communicative problems produces similar neural changes in Senders and Receivers \citep{stolk2013}, and that cerebral coherence between partners appears only in pairs with a shared communicative history \citep{stolk2014}; disrupting the posterior superior temporal sulcus impairs comprehension \citep{stolk2014b} and ventromedial prefrontal lesions alter communicative decisions \citep{stolk2015}. The \emph{conceptual alignment} framework \citep{stolk2016} proposes that communicators maintain a model of what a signal means \emph{to this partner}. This predicts that understanding is partner-specific, consistent with our phase-2 dip when partners change and with reunited pairs doing best in phase 3.

\paragraph{Computational accounts.} Recipient design and intention recognition framed as Bayesian inverse planning are intractable in general and tractable only under constraints \citep{vanrooij2011}. Fitting trial-by-trial TCG data, \citet{blokpoel2012} found Senders' adjustments better explained by perspective taking than by simple heuristics \citep[see also][]{blokpoel2013}. These accounts locate the design work in the Sender, in line with our finding that phase-3 gains come mainly from clearer encoding.

\paragraph{Strategies and individual differences.} Senders typically pause on the goal tile to mark location and convey orientation with a ``wiggle'' to an adjacent tile and back \citep{blokpoel2012}. \citet{volman2012} distinguished coarse from refined strategies; pairs reached about 49\% accuracy on novel and 95\% on familiar problems, and Receiver fluid intelligence predicted performance. Autistic adults produce intelligible signals but struggle when meaning depends on recent shared history \citep{wadge2019}, and adapt less dynamically to partners \citep{koch2026}; five-year-olds already tailor signals to their addressee \citep{stolk2013b}. Our embedding recovers similar dimensions (goal marking, direction marking), but our design separates changes in encoding from changes in decoding.

\paragraph{Conventions in pairs and populations.} The TCG belongs to the wider field of experimental semiotics \citep{galantucci2005,galantucci2011}, alongside movement-only signalling games \citep{scottphillips2009}, coordination games with costly symbols \citep{selten2007} and iterated learning \citep{kirby2008}. Whether conventions spread depends on interaction structure: communities that rotate partners converge on more stable conventions than isolated pairs \citep{garrod1994}; global conventions emerge in well-mixed but not spatially structured networks \citep{centola2015}; partner-specific conventions form quickly in repeated reference games \citep{hawkins2020} and are abstracted into community priors only over many partners \citep{hawkins2023,hawkins2019}; and larger groups or thin channels hamper convergence \citep{boyce2024}. With three phases and few partners, our tournaments sit in the early, dyad-dominated regime these accounts describe, where partner pulls are expected but a group convention is not.

\paragraph{LLMs as participants and similarity judges.} LLMs have been used to simulate survey samples \citep{argyle2023}, replicate classic experiments---with a ``hyper-accuracy distortion'' in newer models \citep{aher2023}---and produce believable social agents \citep{park2023}; a model fine-tuned on trial-level data from 160 experiments predicts held-out participants better than domain models \citep{binz2025}. Critiques question when models can replace participants \citep{dillion2023}, show that they flatten group differences \citep{wang2025}, and question whether apparent success reflects task understanding \citep[preprint]{liu2025}. Our embeddings follow the triplet odd-one-out approach used to map human object concepts \citep{hebart2020,muttenthaler2022}; LLM similarity judgments recover human-like structure across sensory domains \citep{marjieh2024} and object concepts \citep{du2025}. Our finding that LLM-predicted turns change less than real players' turns echoes the hyper-accuracy distortion.
